\documentclass{article}
\usepackage[utf8]{inputenc}
\usepackage{amsmath}
\usepackage{geometry}
\geometry{margin=2cm}
\begin{document}

\section*{Questão 100 — ENEM 2024 — Ciências da Natureza}

\subsection*{Interpretação e Estratégia de Resolução}
A questão pede que se interprete como as linhas de campo elétrico se modificam ao inserir um cilindro metálico oco entre duas placas carregadas com cargas opostas, representado visualmente.

\paragraph{Termos principais:}
\begin{itemize}
  \item \textbf{Linhas de campo elétrico}: Representação gráfica das trajetórias que cargas positivas seguiriam sob ação do campo elétrico.
  \item \textbf{Condutor metálico}: Material que permite movimentação livre de cargas elétricas e impede a passagem de linhas de campo elétrico para seu interior.
\end{itemize}

\paragraph{O que a questão quer:}
Identificar, entre as alternativas, o diagrama que expressa corretamente a modificação das linhas de campo elétrico quando um condutor metálico oco é colocado entre as placas com cargas opostas.

\paragraph{Estratégia:}
Compreender o campo uniforme produzido entre placas carregadas, avaliar o efeito da inserção do condutor metálico oco (que impede penetração das linhas e provoca desvio delas ao redor), comparar as alternativas para selecionar a representação correta desse fenômeno.

\subsection*{Conteúdo e Mini Revisão}
\begin{table}[h!]
\centering
\begin{tabular}{|p{4cm}|p{5cm}|p{6cm}|}
\hline
\textbf{Conceito} & \textbf{Ideia-chave} & \textbf{Aplicação} \\
\hline
Linhas de campo elétrico & Não cruzam e indicam direção do campo & Linhas devem desviar ao redor do condutor sem se cruzar \\
\hline
Condutor em campo elétrico & Linhas não penetram no condutor, se redistribuem ao redor & Cilindro metálico oco desvia linhas, criando região sem linhas no interior \\
\hline
Campo uniforme entre placas & Linhas paralelas, retas e contínuas & Antes do cilindro, linhas são retas; depois desviam ao redor \\
\hline
\end{tabular}
\caption{Revisão conceito de eletrostática pertinente à questão}
\end{table}

\paragraph{Fundamentos teóricos:}
Entre placas metálicas com cargas opostas há campo elétrico uniforme, onde as linhas são retas, paralelas e direcionadas da placa positiva para a negativa. Quando um condutor metálico é inserido, ele impede que as linhas penetrem seu interior, forçando as linhas a contorná-lo. Esse princípio é usado em experimentos para visualização das linhas por meio de materiais indicadores, como farelo de milho.

\subsection*{Resolução e "Pulo do Gato"}
\paragraph{Resolução:}
\begin{enumerate}
  \item O experimento inicial mostra linhas retas e paralelas entre as placas A e B, demonstradas pelo farelo de milho.
  \item Inserindo o cilindro metálico oco, as linhas de campo desviam para contornar o condutor, não penetrando seu interior.
  \item Observando as alternativas, a letra E apresenta esse desvio correto das linhas.
  \item Alternativas A, B, C e D apresentam erros conceituais: penetração no condutor (A), inversão da direção do campo (B), desvios incorretos (C) e cruzamentos de linhas de campo (D).
  \item Portanto, a alternativa correta é \textbf{E}.
\end{enumerate}

\paragraph{Pulo do gato:}
Lembre-se: \textit{linhas de campo elétrico nunca entram em condutores; elas sempre desviam para contornar corpos metálicos, mantendo continuidade fora do condutor}.

\subsection*{Análise das Alternativas e Armadilhas}
\begin{description}
  \item[Alternativa A:] Possível erro por supor que o campo elétrico continua uniforme com linhas atravessando o cilindro; ignora que condutores bloqueiam as linhas.
  \item[Alternativa B:] Ilusão de inversão do campo devido ao condutor, mas isso é incorreto pois o condutor apenas desvia as linhas.
  \item[Alternativa C:] Apresenta desvios incorretos e direções erradas das linhas, indicando má compreensão do fenômeno.
  \item[Alternativa D:] Mostra linhas cruzadas, o que não ocorre em campos elétricos reais; linhas nunca se cruzam.
\end{description}

\subsection*{Código da Questão}
\begin{itemize}
  \item CN\_ELETROSTATICA\_MULTIPLA\_001
\end{itemize}

\bigskip
\noindent \textbf{\large Alternativa correta: E}

\bigskip
\noindent \rule{\textwidth}{0.5pt}

\subsection*{Imagens que acompanham o enunciado e alternativas}

\begin{itemize}
  \item \textit{Imagem do enunciado:} Representação das linhas retas e paralelas do campo uniforme entre as placas A (+) e B (-).
  \item \textit{Imagem alternativa A:} Linhas passando pelo interior do cilindro metálico.
  \item \textit{Imagem alternativa B:} Linhas invertidas em direção, confundindo o campo.
  \item \textit{Imagem alternativa C:} Linhas desviadas com direção incorreta.
  \item \textit{Imagem alternativa D:} Linhas cruzando-se dentro do cilindro.
  \item \textit{Imagem alternativa E:} Linhas desviando corretamente ao redor do cilindro metálico oco, sem ultrapassar seu interior.
\end{itemize}

\end{document}